\section{Nontrivial roots and Gamma amplitudes}\label{cms:sec:roots}
\subsection{The root filtered tree series}
Let $\zeta\ne1$ have exact order $q\geq2$, and put $z=\zeta\ee^{-s}$. In the term of index $r$ of $H_L$, a polylogarithmic singularity occurs precisely when $\zeta^r=1$, or equivalently $q\mid r$. For these terms, the same sign calculation as at the positive root applies. Define
\begin{equation}\label{cms:eq:BqL}
B_{q,L}(s)=\sum_{\substack{1\leq r\leq L\\q\mid r}}
\frac{r^{r-1}}{r!}\,s^{r-1}\ee^{-rs}.
\end{equation}
The corresponding full analytic power series near zero can be expressed as
\begin{equation}\label{cms:eq:Bq}
B_q(s)=\frac1{qs}\sum_{h=0}^{q-1}
 T\bigl(\exp(2\pi\ii h/q)\,s\ee^{-s}\bigr).
\end{equation}
The apparent singularity at $s=0$ is removable. Formula \eqref{cms:eq:Bq} is simply the root-of-unity filter applied to \eqref{cms:eq:tree}. Its first term is
\begin{equation}\label{cms:eq:Bqfirst}
B_q(s)=\frac{q^{q-2}}{(q-1)!}s^{q-1}+O(s^q).
\end{equation}
At any fixed Taylor degree, $B_{q,L}$ agrees with $B_q$ once $L$ is large enough. Its omitted terms have degree at least $L$.

Let $\varepsilon_\zeta=1$ if $\zeta=\pm\ii$ and $0$ otherwise. For a sufficiently small slit neighborhood,
\begin{equation}\label{cms:eq:rootform}
P_L(\zeta\ee^{-s})=s^{\varepsilon_\zeta}U_{\zeta,L}(s)
 \exp\bigl(-B_{q,L}(s)\log s\bigr),
\end{equation}
where $U_{\zeta,L}$ is analytic and nonzero at zero. The factor $s^{\varepsilon_\zeta}$ is precisely the isolated zero of $1+z^2$; it is not absorbed into a branch of its logarithm. The other polylogarithms are analytic at this root, and their exponentials do not vanish.

\subsection{Convergence of the boundary product}
For $\zeta\ne1$ on the unit circle, the series
\[
\sum_{j\geq1}\frac{\zeta^{j+1}}j
\]
converges by Dirichlet's test. For $j\geq2$, the difference between $\log(1+\zeta^{j+1}/j)$ and its linear term is absolutely summable. The logarithm can be chosen from its convergent series because the argument has distance at least $1/2$ from zero. The first factor is dealt with separately. Thus the product has a finite nonzero limit except when its first factor vanishes, which happens at $\zeta=\pm\ii$.

The limit here means the original ordering of factors. At roots of unity, grouping finitely many consecutive residue classes preserves convergence. The same value is approached from inside the disk: the linear logarithmic term obeys Abel's theorem, and the quadratic remainder converges uniformly by its $j^{-2}$ majorant. This also identifies the product limit with the constant term of the finite boundary jet.

\begin{proposition}[Gamma product for a root value]\label{cms:prop:gamma}
If $\zeta$ has exact order $q\geq2$ and $\zeta\ne\pm\ii$, then
\begin{equation}\label{cms:eq:gamma}
F(\zeta)=\prod_{r=1}^{q}
\frac{\Gamma(r/q)}{\Gamma\bigl((r+\zeta^{r+1})/q\bigr)}\ne0.
\end{equation}
\end{proposition}
\begin{proof}
Set $c_r=\zeta^{r+1}$. For the product over $j=q\ell+r$, $0\leq\ell<N$, the Gamma recurrence gives
\begin{align*}
\prod_{\ell=0}^{N-1}\left(1+\frac{c_r}{q\ell+r}\right)
&=\prod_{\ell=0}^{N-1}\frac{\ell+(r+c_r)/q}{\ell+r/q}\\
&=\frac{\Gamma\bigl(N+(r+c_r)/q\bigr)}{\Gamma(N+r/q)}
\frac{\Gamma(r/q)}{\Gamma\bigl((r+c_r)/q\bigr)}.
\end{align*}
The ratio depending on $N$ is $N^{c_r/q}(1+O(N^{-1}))$. Multiplication over $r=1,\ldots,q$ cancels all powers of $N$, since
\[
\sum_{r=1}^{q}c_r=\sum_{r=1}^{q}\zeta^{r+1}=0.
\]
Taking $N\to\infty$ proves \eqref{cms:eq:gamma}.

A factor of the original product can vanish only if $j=1$ and $\zeta^2=-1$. Equivalently, in the Gamma expression only $r=1$, $c_1=-1$ can make a denominator argument zero or a negative integer: for $r\geq2$, its real part is at least $(r-1)/q>0$. The excluded roots are therefore the only zeros. This also follows directly from convergence of the logarithmic remainder.
\end{proof}

\subsection{The leading singularity and the natural boundary}
For $\zeta\ne1,\pm\ii$ of exact order $q$, let $t=1-z/\zeta$. We have $s=-\log(1-t)=t+O(t^2)$. Equations \eqref{cms:eq:rootform} and \eqref{cms:eq:Bqfirst}, multiplied by the Taylor jet of $Q_L$, show that the first nonanalytic term is
\begin{equation}\label{cms:eq:rootleading}
-F(\zeta)\frac{q^{q-2}}{(q-1)!}\,
 t^{q-1}\log t.
\end{equation}
Its coefficient is nonzero. Choosing $L$ so that the smooth tail has more than $q$ derivatives rules out cancellation of the logarithmic divergence by the remainder. In particular, $F$ does not extend holomorphically through this root.

\begin{corollary}\label{cms:cor:natural}
The unit circle is a natural boundary of $F$.
\end{corollary}
\begin{proof}
Roots of unity other than $1$ and $\pm\ii$ are dense on the circle, and each is singular by \eqref{cms:eq:rootleading}. Analytic continuation through any open arc would contradict one of those singularities. The radius of convergence is exactly $1$, already ensured by the pole at $1$.
\end{proof}

For later reference, the generic leading coefficient from a single order-$q$ root, expressed in the coefficient index $m$, is
\begin{equation}\label{cms:eq:genericcoefficient}
(-1)^{q+1}q^{q-2}F(\zeta)\zeta^{-m}m^{-q}.
\end{equation}
This follows from the transfer formula in Section~\ref{cms:sec:transfer}. It is the leading contribution of that root, not the whole term of order $m^{-q}$, since roots of smaller order also have higher corrections.

\subsection{The zeros at the order four roots}
Define the product with its first factor removed by
\[
\widehat F(z)=\frac{F(z)}{1+z^2}
 =\prod_{j\geq2}\left(1+\frac{z^{j+1}}j\right).
\]
The same convergence proof gives nonzero boundary values at $\pm\ii$. Grouping the factors at $\ii$ gives
\begin{equation}\label{cms:eq:Qi}
Q_{\ii}:=\widehat F(\ii)
=\frac{\Gamma(5/4)\Gamma(1/2)\Gamma(3/4)}
 {\Gamma((2-\ii)/4)\Gamma((4+\ii)/4)}\ne0.
\end{equation}
For example, the residue class $j\equiv1\pmod4$ starts at $j=5$ because $j=1$ has been removed; its Gamma ratio is $\Gamma(5/4)/\Gamma(1)$. The other three residue classes give the remaining factors in \eqref{cms:eq:Qi}.

At $z=\ii\ee^{-s}$,
\[
1+z^2=1-\ee^{-2s}=2s+O(s^2),\qquad
B_4(s)=\frac83s^3+O(s^4).
\]
Hence the first nonanalytic term of $F$ is
\begin{equation}\label{cms:eq:izero}
-\frac{16}{3}Q_{\ii}s^4\log s
=-\frac{16}{3}Q_{\ii}t^4\log t+\text{higher nonanalytic terms}.
\end{equation}
Thus the order-four roots first affect coefficients at order $m^{-5}$, rather than $m^{-4}$. Indeed their conjugate pair contributes at its first order
\begin{equation}\label{cms:eq:ipair}
256\Rea(Q_{\ii}\ii^{-m})m^{-5}.
\end{equation}
This is only the contribution of this pair at that order. It is not a claimed complete formula for $P_5$. The zero delays the logarithmic singularity; it does not remove it. In particular $\pm\ii$ are singular as well, although they were unnecessary for the density argument.
